\chapter{HUD and Menu System}

\section{Retained HUD Workflow}

DemiEngine uses declarative \filepath{*.hud.json} files for authored UI. Lua
may generate genuinely dynamic nodes through \code{Hud.create}; both paths use
the same retained tree, layout, input, and accessibility systems.

\section{Declarative HUD Elements}

HUD files support retained layout and control nodes, including:

\begin{longtable}{>{\ttfamily}p{0.20\textwidth}p{0.68\textwidth}}
panel & A colored container with position, size, anchors, and child nodes. \\
label & Text with font size, color, alignment, and wrapping. \\
button & A focusable control with label, style, and typed action. \\
\end{longtable}

\section{Canvas Size}

Each HUD file can define its own coordinate system:

\begin{lstlisting}[language=json]
"canvas_size": [1920.0, 1080.0]
\end{lstlisting}

All positions and sizes in that HUD file are authored against this canvas. At runtime, the renderer scales the HUD to the actual window size. HUD click detection uses the same conversion, so fullscreen and resized windows remain interactive.

\section{Groups}

Groups are ordinary structural parent nodes. Toggling the parent controls its
entire subtree:

\begin{lstlisting}[language={[5.2]Lua}]
Hud.set_visible("menu_main", true)
Hud.set_visible("menu_options", false)
\end{lstlisting}

The \filepath{minimal\_2d} menu uses these groups:

\begin{itemize}
  \item \code{menu\_base};
  \item \code{menu\_main};
  \item \code{menu\_options};
  \item \code{menu\_sound};
  \item \code{menu\_video};
  \item \code{menu\_dropdown};
  \item \code{game\_hud}.
\end{itemize}

\section{Button Events}

Buttons can declare an action annotation:

\begin{lstlisting}[language=json]
{
  "type": "button",
  "id": "menu_button_start",
  "label": "START GAME",
  "action": "menu_button_start"
}
\end{lstlisting}

The runtime dispatches the action to Lua functions annotated with \code{@HandleAction}:

\begin{lstlisting}[language=lua]
-- @HandleAction("menu_button_start")
function Menu.start_game(event)
  Scene.load("scene://game/level_1")
end
\end{lstlisting}

Typed UI events are also available through \code{on\_ui\_event}, specialized
callbacks such as \code{on\_ui\_submit}, and matching event-bus channels.

\section{Event Annotations}

Lua functions can subscribe to game events with \code{@OnEvent} instead of calling \code{Events.subscribe} manually:

\begin{lstlisting}[language=lua]
-- @OnEvent("player_died")
function GameEvents.show_game_over(event)
  Hud.set_text("game_over_points", "POINTS: " .. tostring(event.points))
end
\end{lstlisting}

For reusable modules loaded with \code{require}, list the module in \code{scripting.modules} in the project file so the runtime can register its annotations.
